% Week 2 Session A — Chain-of-Thought prompting practice
% All lab reports in this series include an abstract.
%
% Build: pdflatex Week2_SessionA_Report.tex from Lab-reports/
% Same folder: week2_session_a_opencode.png, week2_session_a_opencode2.png,
%   week2_session_a_opencode3.jpg, week2_session_a_terminal.png,
%   prompt_log_week2_session_a.md
%
\documentclass[11pt,a4paper]{article}

\usepackage[margin=2.5cm]{geometry}
\usepackage{graphicx}
% Tall OpenCode screenshots: limit height so each figure fits on one page
\newcommand{\opencodefig}[1]{%
  \includegraphics[width=\linewidth,height=0.68\textheight,keepaspectratio]{#1}%
}
\newcommand{\terminalfig}[1]{%
  \includegraphics[width=0.92\linewidth,height=0.38\textheight,keepaspectratio]{#1}%
}
\usepackage{booktabs}
\usepackage{xurl}
\usepackage{hyperref}
\usepackage{parskip}
\usepackage{enumitem}
\usepackage{xcolor}
\usepackage{fvextra}
\usepackage[most]{tcolorbox}

\newcommand{\studentname}{Mahmudul Hasan}
\newcommand{\studentid}{4125999049}
\newcommand{\reportdate}{May 17, 2026}

\makeatletter
\g@addto@macro\UrlBreaks{\do\ }
\makeatother

\newcommand{\LabCodeRoot}{.}

\title{Lab Report: Week 2 Session A\\Chain-of-Thought Prompting Practice}
\author{\studentname \quad (Student ID: \studentid)}
\date{\reportdate}

\begin{document}
\maketitle

\begin{abstract}
This report documents Week~2 Session~A: advanced \textbf{chain-of-thought (CoT)} prompting with OpenCode Desktop on (1)~a five-step rainfall monitoring system design and (2)~line-by-line debugging of a self-prepared SCS-CN script. A shorter design prompt produced a large implementation stack; CoT debugging identified four formula errors. Local runs gave $Q \approx 0.002$~mm (buggy) and $Q = 43.6$~mm (fixed) for $P=80$~mm, $\mathrm{CN}=85$, contradicting an incorrect closing estimate in the AI chat. Evidence: four screenshots and \path{prompt_log.md} in Appendix~\ref{app:promptlog}.
\end{abstract}

\section{Introduction}
Week~2 Session~A extends Week~1 by using CoT for \emph{complex design} and \emph{debugging}, not only single-formula calculations. Work was completed in \path{ai_water_lab/week2_session_a/} with OpenCode Desktop (OpenCode Zen). Exercise~1 compares a structured five-step design prompt with a one-line urban flood-warning prompt. Exercise~2 uses a student-authored \path{buggy_rainfall.py} (instructor directed the class to prepare their own buggy code) and CoT-assisted repair verified in the terminal.

\section{Objectives}
\begin{itemize}[noitemsep]
  \item Practice CoT for multi-step hydrological system design.
  \item Compare CoT vs.\ shorter prompting on the same design goal.
  \item Debug Python using CoT (line-by-line explanation and fixes).
  \item Document outcomes in a prompt log and personal prompting strategy.
\end{itemize}

\section{Methods}

\subsection{Exercise 1: Rainfall monitoring design}
The CoT prompt asked for reasoning on: data requirements, real-time API, storage structure, flood alert thresholds, and visualization (with code suggestions where appropriate). A follow-up used: \emph{Design a rainfall monitoring system for urban flood warning.}

\subsection{Exercise 2: Debugging with CoT}
A short SCS-CN script with intentional errors was written for practice. Before AI assistance, four bugs were hypothesized (wrong $S$ formula, wrong $I_a$, missing $P \le I_a$ branch, wrong denominator). OpenCode received:

\emph{This code has bugs. Walk through it line by line, explain what each part should do, identify what's wrong, and suggest fixes.}

Fixes were implemented in \path{rainfall_fixed.py} and tested with \texttt{python3}.

\section{Results}

\subsection{Exercise 1: CoT design response}
OpenCode recommended core and derived rainfall fields, Open-Meteo for prototyping, TimescaleDB hypertables, tiered WMO/NWS-style alert levels, and a Plotly/Grafana visualization path with an API$\rightarrow$fetcher$\rightarrow$database$\rightarrow$alert pipeline. Figure~\ref{fig:cot-design} shows the CoT exchange.

\begin{figure}[htbp]
  \centering
  \opencodefig{week2_session_a_opencode.png}
  \caption{OpenCode: five-step CoT prompt for rainfall monitoring system design (Exercise~1).}
  \label{fig:cot-design}
\end{figure}

\subsection{Exercise 1: Short prompt response}
The short prompt generated a multi-file urban flood stack (\path{docker-compose.yml}, FastAPI, zone-based flood model, alert router, etc.) rather than a stepwise design narrative. Figure~\ref{fig:short-design} documents this response.

\begin{figure}[htbp]
  \centering
  \opencodefig{week2_session_a_opencode2.png}
  \caption{OpenCode: short prompt \emph{Design a rainfall monitoring system for urban flood warning} (Exercise~1 comparison).}
  \label{fig:short-design}
\end{figure}

\subsection{Exercise 2: CoT debugging}
OpenCode identified the same four bugs as the pre-AI list and supplied a corrected \texttt{calculate\_runoff} function. Figure~\ref{fig:debug} shows the session.

\begin{figure}[htbp]
  \centering
  \opencodefig{week2_session_a_opencode3.jpg}
  \caption{OpenCode: CoT line-by-line debug of \texttt{buggy\_rainfall.py} (Exercise~2).}
  \label{fig:debug}
\end{figure}

\subsection{Exercise 2: Terminal verification}
Local execution confirmed the failure mode and fix (Figure~\ref{fig:terminal}):

\begin{itemize}[noitemsep]
  \item \path{buggy_rainfall.py}: $Q \approx 0.0019$~mm (unrealistic for $P=80$~mm).
  \item \path{rainfall_fixed.py}: $Q = 43.6$~mm (consistent with Week~1 Session~B SCS-CN check).
\end{itemize}

\begin{figure}[htbp]
  \centering
  \terminalfig{week2_session_a_terminal.png}
  \caption{Terminal: before/after fix for $P=80$~mm, $\mathrm{CN}=85$.}
  \label{fig:terminal}
\end{figure}

Table~\ref{tab:prompts} summarizes Exercise~1 prompting styles.

\begin{table}[htbp]
  \centering
  \caption{CoT vs.\ short prompt for rainfall monitoring design.}
  \label{tab:prompts}
  \begin{tabular}{@{}p{0.22\linewidth}p{0.34\linewidth}p{0.34\linewidth}@{}}
    \toprule
    Aspect & CoT (5 steps) & Short (one line) \\
    \midrule
    Output & Design rationale per step & Many code files + Docker stack \\
    Audit trail & Strong & Weaker unless you inspect each file \\
    Best use & Planning, reports, grading & Rapid prototype after design is fixed \\
    \bottomrule
  \end{tabular}
\end{table}

\section{Analysis and validation}

\textbf{Design (Exercise~1):} Recommendations are physically plausible for a conceptual monitoring system. Thresholds and APIs must still be checked against course labs (e.g.\ OpenWeatherMap in Week~1 Session~A; 20~mm/h alert in Week~5 Lab~1). Open-Meteo is reasonable for keyless prototyping.

\textbf{Debugging (Exercise~2):} All four intentional bugs were found by both the author and OpenCode. The AI's closing remark that the fully buggy script might yield $\approx 48.9$~mm on this input was \textbf{not} supported by the terminal ($\approx 0.002$~mm), because an inflated $S$ from multiplication dominates the denominator---an example of why local verification is required.

\textbf{Prompting strategy:} CoT improved traceability for design and debugging. Short prompts can accelerate coding but may skip explicit reasoning.

\section{Discussion}
\begin{enumerate}[noitemsep]
  \item CoT produced clearer structure for the five design questions; the short prompt jumped to implementation.
  \item For hydrology coursework, CoT debugging plus terminal tests caught an AI numeric slip that line-by-step reasoning alone might miss.
  \item Limits of AI output include unverified production claims (Docker stack) and occasional wrong numeric summaries.
  \item Next time: specify deliverable type in the prompt (design memo vs.\ code only) and always run corrected scripts locally.
\end{enumerate}

\section{Problems encountered}
No instructor buggy file was distributed; the class was told to create their own. The main extra step was renaming screenshots and copying \path{prompt_log.md} for the report appendix.

\section{Lessons learned}
Numbered CoT prompts align AI answers with lab rubrics. Comparing prompt styles on the same topic clarifies when to ask for reasoning first and code second. A prompt log plus screenshots and terminal output provides defensible evidence for grading.

\section{Conclusion}
Week~2 Session~A objectives were met: CoT system design, CoT vs.\ short comparison, CoT debugging of SCS-CN code, and documentation in \path{prompt_log.md}. Course submission: \path{prompt_log.md}, \path{buggy_rainfall.py}, \path{rainfall_fixed.py}, and screenshots as specified in the lab guide.

\appendix

\section{Submitted prompt log: \texttt{prompt\_log.md}}
\label{app:promptlog}
\begin{tcolorbox}[
  enhanced,
  colback=gray!6,
  colframe=black!55,
  title={\texttt{prompt\_log.md} \mdseries (full file)},
  fonttitle=\bfseries\ttfamily\footnotesize,
  arc=1.5mm,
  boxrule=0.7pt,
  breakable,
  width=\linewidth,
  before skip=0.8em,
  after skip=0.8em,
]
\VerbatimInput[
  fontsize=\footnotesize,
  breaklines=true,
]{\LabCodeRoot/prompt_log_week2_session_a_latex.md}
\end{tcolorbox}

\end{document}
